\chapter{Primary-Source Guide and Theorem Crosswalk}
\label{app:proetale-references}

This appendix identifies the primary source of each deep theorem and records
the convention under which it is used.  Expository notes are valuable route
maps, but when a shift, twist, hypothesis, or logical dependency is delicate,
the primary theorem listed here governs this text.

\section{Bhatt--Scholze: the pro-etale model}

The primary coefficient source is Bhatt--Scholze,
\href{https://people.mpim-bonn.mpg.de/scholze/proetale.pdf}
{\emph{The pro-etale topology for schemes}}, Astérisque 369 (2015).

\begin{longtable}{
 >{\raggedright\arraybackslash}p{0.28\textwidth}
 >{\raggedright\arraybackslash}p{0.20\textwidth}
 >{\raggedright\arraybackslash}p{0.42\textwidth}}
\toprule
Topic & Primary location & Use in these notes\\
\midrule
\endhead
Weakly etale and pro-etale &
Definition 1.2; Theorem 1.3 &
Definition of \(X_{\proet}\), and comparison of weakly etale with
ind-etale maps.\\

Weakly contractible basis &
Definition 1.4; Theorem 1.5; Theorem 1.8 &
\cref{thm:proetale-wcontractible-cover}; exact evaluation and control of
inverse limits.\\

Replete topoi and inverse limits &
Propositions 3.1.10--3.1.11; Section 3.3 &
Surjective systems satisfy
\(\varprojlim\simeq\Rlim\); countable inverse limit has cohomological
dimension one; left completeness.\\

Derived completion in a replete topos &
Proposition 3.5.1; Lemma 3.5.7 &
For a replete topos, a noetherian ring, and an ideal \(\mathfrak m\), derived
completeness is reconstruction from the reductions.  Complete complexes are
compared with compatible systems in the bounded-above range; the unbounded
comparison additionally assumes that \(\mathfrak m\) is regular.\\

Comparison with the etale site &
Corollaries 5.1.5--5.1.6 &
Discrete local systems agree, and bounded-below torsion complexes have the
same cohomology after pullback by \(\nu\).\\

Complete constructible objects &
Definition 6.5.1; Lemmas 6.5.3--6.5.9 &
\cref{thm:proetale-complete-system-equivalence}; derived completeness,
finite reduction, tensor closure, and locally closed functors.\\

Stratified characterization &
Proposition 6.6.11 &
\cref{thm:proetale-constructible-characterization}: constructible exactly
means locally perfect along a finite stratification when \(X\) is
topologically noetherian.\\

Pushforward and base change &
Theorem 6.7.1; Lemmas 6.7.2, 6.7.4--6.7.5 &
Constructibility of \(Rf_*\), smooth base change, and proper base change
under the stated finiteness hypotheses.\\

Exceptional direct image &
Definition 6.7.6; Lemma 6.7.7; Lemma 6.7.10 &
Compactification definition of \(Rf_!\), independence, constructibility,
and base change.\\

Failure of the naive derived \(f_!\) &
Remarks 6.7.8--6.7.9 &
\Cref{warn:dictionary-fshriek}; prevents an incorrect abelian-derived
construction.\\

Tensor, internal Hom, projection &
Lemmas 6.7.12--6.7.14 &
Closed monoidal constructible category and projection formula.\\

Exceptional inverse image &
Lemmas 6.7.18--6.7.19 &
Construction of \(Rf^!\) from compatible finite-level adjoints.\\

Dualizing complexes &
Lemma 6.7.20 &
Integral Verdier duality, built from the finite-level dualizing complexes.\\

Integral structures for lisse \(E\)-sheaves &
Lemma 6.8.13 &
On a topologically noetherian scheme, every lisse \(E\)-sheaf admits a
constructible \(\mathcal O_E\)-model.  The model need not be lisse globally;
the proof first obtains lattices after a finite stratification.\\

\(\mathcal O_E,E,\Qbarl\)-coefficients &
Definitions 6.8.1--6.8.3; Proposition 6.8.14; Remark 6.8.15 &
\cref{thm:proetale-coefficient-tower,%
thm:proetale-six-functor-package}; integral lattices, inversion of the
uniformizer, finite coefficient-field descent, and six operations.\\

Pro-etale fundamental group &
Theorem 1.10; Section 7 &
Continuous representations on nonnormal schemes; warning against replacing
\(\pi_1^{\proet}\) everywhere by its profinite completion.\\
\bottomrule
\end{longtable}

The hypotheses of Section 6.7 are part of every use: complete noetherian
local coefficients with finite residue field of characteristic \(\ell\);
qcqs and finitely presented maps; and either properness or the relevant
quasi-excellent, \(\ell\)-coprime hypothesis.  Finite-type schemes over
\(\Fq\), with \(\ell\nmid q\), lie in this range.

\section{Modern constructibility by dualizability}

Hemo--Richarz--Scholbach,
\href{https://arxiv.org/abs/2305.18131}
{\emph{Constructible sheaves on schemes}}, Advances in Mathematics 429
(2023), gives the coefficient-flexible categorical formulation.

\begin{longtable}{
 >{\raggedright\arraybackslash}p{0.29\textwidth}
 >{\raggedright\arraybackslash}p{0.20\textwidth}
 >{\raggedright\arraybackslash}p{0.41\textwidth}}
\toprule
Topic & Location & Formalization role\\
\midrule
\endhead
Lisse as dualizable &
Definition 1.1 &
The reusable predicate
\(\operatorname{Lisse}(K)\leftrightarrow
 \operatorname{IsDualizable}(K)\).\\

Weakly contractible calculation &
Lemma 1.2, referring to the detailed theorem in Section 4 &
On a weakly contractible affine, lisse objects are perfect coefficient
complexes.\\

Pro-etale descent &
Lemma 1.3 and the detailed descent results &
Lisse and constructible categories form hypersheaves; choices of
trivializing covers do not enter public objects.\\

Sequential coefficient limits &
Proposition 5.1 &
If \(\Lambda=\varprojlim_n\Lambda_n\) is sequential, with surjective
transitions having locally nilpotent kernels, then lisse and constructible
categories are the corresponding inverse limits.\\

Filtered coefficient colimits &
Proposition 5.2 &
For qcqs \(X\) and \(\Lambda=\varinjlim_i\Lambda_i\), lisse and constructible
categories are the corresponding filtered colimits.\\

Localization of coefficients &
Proposition 5.5; Remark 5.7 &
For qcqs \(X\), localization along a multiplicative subset gives fully
faithful functors on lisse and constructible categories.  Essential
surjectivity is not automatic for lisse objects, reflecting the distinction
between the etale and pro-etale fundamental groups.\\

Natural \(t\)-structure &
Theorem 6.2 &
The coefficient ring must be \(t\)-admissible.  Lisse objects additionally
require that \(X\) have Zariski locally finitely many irreducible components;
constructible objects require this for every constructible subset of \(X\).
These hypotheses are part of the interface.\\

Classical adic comparison &
Theorem 7.7 &
For qcqs \(X\), the pro-etale constructible categories with
\(\mathcal O_E/\mathfrak m^n\), \(\mathcal O_E\), and \(E\) coefficients
agree with the corresponding classical finite-Tor and adic categories.\\

Algebraic coefficient closure &
Lemma 7.9 &
For qcqs \(X\), filtered colimits over finite extensions
\(E/\Ql\) recover the constructible categories with \(\overline{\mathbf
Z}_\ell\) and \(\Qbarl\) coefficients.\\
\bottomrule
\end{longtable}

The notes use Bhatt--Scholze's concrete complete-DVR model for the proof
spine and the dualizability formulation as the reusable public abstraction.

\section{Finite-level six operations, purity, and trace}

The pro-etale formalism lifts finite-level geometric theorems; it does not
replace their proofs.

\begin{longtable}{
 >{\raggedright\arraybackslash}p{0.25\textwidth}
 >{\raggedright\arraybackslash}p{0.27\textwidth}
 >{\raggedright\arraybackslash}p{0.38\textwidth}}
\toprule
Result & Primary source & Use\\
\midrule
\endhead
Etale topoi and torsion cohomology &
SGA 4 \cite{SGA4}, Exposes V--XII &
Finite coefficient derived categories, stalks, cohomological dimension,
and descent.\\

Finiteness and base change &
SGA 4 \cite{SGA4}, Exposes XIV--XVII; SGA \(4\frac12\)
\cite{SGA4half} &
Constructibility and finite amplitude, proper/smooth base change, and
compact support.\\

Exceptional operations &
SGA 4 \cite{SGA4}, Expose XVII &
Compactification, \(Rf_!\), composition, and base change at finite level.\\

Duality and quasi-excellent refinements &
\href{https://www.numdam.org/item/AST_2014__363-364__R1_0/}
{Illusie--Laszlo--Orgogozo, \emph{Travaux de Gabber}},
especially Expose XVII &
Finite-level \(Rf^!\), dualizing complexes, and compatibility in the range
used by Bhatt--Scholze 6.7.18--6.7.20.\\

Absolute cohomological purity &
Gabber's absolute purity theorem, as presented in the same volumes &
Costalk normalization, smooth purity, and Poincare duality.\\

Lefschetz--Verdier trace &
SGA 5 \cite{SGA5}; SGA \(4\frac12\) \cite{SGA4half} &
\Cref{ch:trace-formula}; local terms, Frobenius fixed points, and the
cohomological zeta formula.\\
\bottomrule
\end{longtable}

\section{BBD: perverse sheaves and weights}

The categorical and weight-theoretic source is
Beilinson--Bernstein--Deligne, \emph{Faisceaux pervers} \cite{BBD}.

\begin{longtable}{
 >{\raggedright\arraybackslash}p{0.29\textwidth}
 >{\raggedright\arraybackslash}p{0.19\textwidth}
 >{\raggedright\arraybackslash}p{0.42\textwidth}}
\toprule
Result & BBD location & Use here\\
\midrule
\endhead
Gluing and middle perversity &
Sections 1.4, 2.1--2.2 &
\cref{def:middle-perversity,thm:middle-t-structure}.\\

Integral \(p\) and \(p^+\) &
Section 3.3 &
\cref{thm:p-pplus-tilt,thm:integral-verdier-exchange};
duality exchanges the two hearts.\\

Affine and quasi-finite exactness &
Theorem 4.1.1 and Corollaries 4.1.2--4.1.3 &
\cref{prop:basic-perverse-exactness,thm:perverse-artin-engine}.\\

Smooth and product exactness &
4.2.4 and 4.2.7--4.2.8 &
Dimension-shifted smooth pullback and external products.\\

Finite length and simple objects &
Theorem 4.3.1 &
\cref{thm:perverse-finite-length,thm:simple-perverse-classification}.\\

Geometric positive-degree Ext vanishing &
5.1.15(iii) &
Splitting arithmetic extensions after geometric base change; no claim of
Frobenius semisimplicity.\\

Intermediate extension and weights &
Corollary 5.3.2 &
Purity of \(j_{!*}\), used only after the upper direct-image theorem.\\

Pure perverse semisimplicity &
Theorem 5.3.8 &
Geometric semisimplicity of the rational perverse object.\\

Ordinary/perverse weight criterion &
Theorem 5.4.1 and Corollary 5.4.4 &
Conversion between derived weights and weights on
\(\pH^a\).\\

Splitting of pure complexes &
Theorem 5.4.5 &
Noncanonical splitting after geometric base change.\\

Decomposition theorem &
Theorem 6.2.5 &
\cref{thm:decomposition-engine,thm:decomposition-finite-field};
a consequence of purity, not an input to it.\\
\bottomrule
\end{longtable}

For an accessible discussion of the integral failure of self-duality and
the \(p/p^+\) pair, see Juteau--Mautner--Williamson,
\href{https://arxiv.org/abs/0901.3322}
{\emph{Perverse sheaves and modular representation theory}}, Section 2.3,
which points back to BBD 3.3.

\section{Deligne: the arithmetic weight theorem}

Deligne's \emph{La conjecture de Weil II} \cite{DeligneWeilII} supplies the
arithmetic estimates and their exact noncircular order.

\begin{longtable}{
 >{\raggedright\arraybackslash}p{0.28\textwidth}
 >{\raggedright\arraybackslash}p{0.18\textwidth}
 >{\raggedright\arraybackslash}p{0.44\textwidth}}
\toprule
Result & Location & Role\\
\midrule
\endhead
Weights and determinant weights &
Sections 1.2--1.3 &
Definitions, tensor/dual/twist bookkeeping, and rank-one inputs.\\

Real-sheaf criterion &
Theorem 1.5.1 &
Mixedness mechanism used before the general direct-image theorem.\\

Local monodromy weights &
Theorem 1.8.4 &
Boundary terms in the local Fourier and curve arguments.\\

Curve theorem &
Theorem 3.2.3 &
\cref{thm:laumon-curve-engine}; reconstructed through perverse Fourier
transform.\\

General direct-image upper bound &
Theorem 3.3.1 &
\cref{thm:perverse-direct-image-weight-bound,%
thm:deligne-laumon-direct-image}.\\

Cohomological upper and lower bounds &
Corollaries 3.3.4--3.3.6 &
Corollary 3.3.4 says that \(H_c^i\) sends weights \(\leq n\) to weights
\(\leq n+i\).  Corollary 3.3.5 gives the dual lower bound for \(H^i\) on a
smooth scheme with lisse coefficients.  Corollary 3.3.6 says that, for a
pure lisse sheaf of weight \(n\) on a smooth scheme, the image of
\(H_c^i\to H^i\) is pure of weight \(n+i\).  Proper purity follows in the
proper case because this map is an isomorphism; it is not the literal
statement of all three corollaries.\\

Smooth-proper consequence &
Corollary 3.3.9 &
\cref{thm:weil-rh-perverse-engine,%
cor:perverse-smooth-proper-purity}.\\
\bottomrule
\end{longtable}

\section{Laumon: Fourier transform as a perverse equivalence}

The Fourier route is from Laumon,
\emph{Transformation de Fourier, constantes d'equations fonctionnelles et
conjecture de Weil} \cite{LaumonFourier}.

\begin{longtable}{
 >{\raggedright\arraybackslash}p{0.28\textwidth}
 >{\raggedright\arraybackslash}p{0.20\textwidth}
 >{\raggedright\arraybackslash}p{0.42\textwidth}}
\toprule
Result & Location & Normalization/use\\
\midrule
\endhead
Definition &
1.2.1.1 &
\[
 \FT_\psi(K)
  =Rp_{2!}
   \bigl(p_1^*K\otimes\cL_\psi(\langle x,\xi\rangle)\bigr)[r].
\]\\

Inversion &
Theorem 1.2.2.1 &
Inverse character gives the inverse transform with the prescribed Tate
twist; using the same character inserts pullback by the sign map.\\

Support-forgetting &
Theorem 1.3.1.1 &
The Fourier-kernel map \(Rp_!\to Rp_*\) is an isomorphism.\\

Duality &
Corollary 1.3.2.2 &
Verdier duality replaces \(\psi\) by \(\psi^{-1}\) and contributes the
displayed twist.\\

Perverse \(t\)-exactness &
Theorem 1.3.2.3 and Corollary 1.3.2.4 &
\cref{thm:fourier-perverse-equivalence-engine,%
thm:fourier-perverse-texact-detailed}; exact equivalence of perverse hearts.\\

Local Fourier transforms &
Sections 2.3--2.5 &
Boundary monodromy and conductor contributions.\\

Curve-weight proof &
Sections 4.1--4.4 &
The perverse curve engine and the reduction to the global direct-image
theorem.\\
\bottomrule
\end{longtable}

The additive character is allowed to enlarge the finite coefficient field
\(E\).  The Fourier kernel has an \(\mathcal O_E\)-lattice; it is not
asserted to be defined over bare \(\Zl\).

\section{Classical Weil assembly}

The trace-to-zeta and duality-to-functional-equation steps follow the
classical cohomological architecture in SGA 5 \cite{SGA5}, Deligne
\cite{DeligneWeilI,DeligneWeilII}, and standard references such as
Freitag--Kiehl \cite{FreitagKiehl}.  In these notes, however, their algebraic
content is proved independently in the algebraic-core chapter:

\begin{center}
\begin{tabularx}{0.94\textwidth}{>{\raggedright\arraybackslash}p{0.32\textwidth}X}
\toprule
Geometric input & Geometry-free output\\
\midrule
Trace formula &
Alternating determinant expression for \(Z(X,T)\).\\
Perfect Poincare pairing &
Reciprocal characteristic polynomials and functional equation.\\
Universal-realization purity &
Separation of factors of distinct weights.\\
Integral rational series &
Integral individual factors by the Fatou lattice and Gauss lemma.\\
\bottomrule
\end{tabularx}
\end{center}

\section{Convention crosswalk}

\begin{longtable}{
 >{\raggedright\arraybackslash}p{0.30\textwidth}
 >{\raggedright\arraybackslash}p{0.29\textwidth}
 >{\raggedright\arraybackslash}p{0.31\textwidth}}
\toprule
Convention in these notes & Consequence & Conversion warning\\
\midrule
\endhead
Geometric Frobenius &
\(F\mid E(1)=q^{-1}\). &
Arithmetic Frobenius inverts every eigenvalue.\\

Derived weight convention &
\(\cH^aK\) has weight \(w+a\) when \(K\) is pure of weight \(w\). &
Thus \(K[n]\) is pure of weight \(w+n\).\\

Perverse-normalized IC &
\(\IC_X(\cL)=j_{!*}(\cL[d])\). &
An unshifted intersection complex differs by \([-d]\).\\

Integral perversity &
Verdier duality sends \(p\) to \(p^+\). &
Self-duality is asserted only after rationalization.\\

Fourier normalization &
The rank-\(r\) transform includes \([r]\). &
Without this shift it is not perverse \(t\)-exact.\\

Fourier duality &
\(\psi\) changes to \(\psi^{-1}\). &
Using the same character also changes the sign-map convention.\\

Frobenius polynomial &
\(P_i(T)=\det(1-TF\mid H^i)\). &
Polynomial roots are reciprocals of Frobenius eigenvalues.\\

Comparison categories &
Classical and pro-etale objects correspond by an equivalence. &
The equivalence and its compatibilities are not definitional equalities.\\
\bottomrule
\end{longtable}

These choices determine every shift, twist, exponent, and absolute value in
the terminal theorem.

\section{Expository route maps}

Kiehl--Weissauer \cite{KiehlWeissauer} remains the closest book-length guide
to the perverse Fourier proof.  The Stanford Weil II seminar
\cite{StanfordWeilIISeminar}, especially Lectures 18--24
\cite{StanfordL18,StanfordL19,StanfordL22,StanfordL23,StanfordL24}, gives a
useful reading sequence.  These sources predate or do not center the direct
pro-etale coefficient model; Bhatt--Scholze governs that replacement.
